\documentclass[aspectratio=169]{beamer}
\input{header}
\coursecode{JKSSB}

\title{Current-Paper Depth Audit \& Technical-Post Routing}
\subtitle{Lesson 66 --- Staying Current and Routing by Post}
\author{JKSSB Computer Awareness}
\date{\today}

\begin{document}
\frame{\titlepage}

\begin{frame}{What This Lesson Is (and Is Not)}
  \begin{itemize}
    \item This is a \key{meta lesson}: a method for staying current, not a new topic.
    \item It teaches a repeatable \key{depth-audit workflow}: check the latest paper, tag each item, route study by post.
    \item It does not add new facts about hardware, Office, or networks.
    \item Everything here applies to \key{every earlier lesson} in the course.
  \end{itemize}
  \begin{block}{The one idea}
    Study to the depth your post actually tests; verify every dispute against the official paper and key.
  \end{block}
\end{frame}

\begin{frame}{Why Papers Drift (and Why Notes Go Stale)}
  \begin{itemize}
    \item JKSSB issues a fresh \key{notification} for each post, with its own scope and marks.
    \item Each cycle brings a new \key{question paper} and a \key{provisional key}, then a \key{final key}.
    \item Software and protocols change: Office menus move, protocol versions advance.
    \item A note written last year is \key{stale} until it is checked against the latest Tier-1 sources.
  \end{itemize}
  \begin{alertblock}{Takeaway}
    Never freeze preparation; re-audit notes against the newest paper and final key each cycle.
  \end{alertblock}
\end{frame}

\begin{frame}{The Revision Workflow in Five Steps}
  \begin{itemize}
    \item \key{Step 1:} collect the newest Tier-1 set: notification, paper, provisional key, final key.
    \item \key{Step 2:} list every computer item and \key{tag} it by depth (next frame).
    \item \key{Step 3:} compare the \key{provisional key} with the \key{final key} for changed answers.
    \item \key{Step 4:} mark \key{version-sensitive} claims (Office version, protocol facts) for re-checking.
    \item \key{Step 5:} route: study your post's tags deeply, read other tags lightly.
  \end{itemize}
  \muted{Run this workflow once per recruitment cycle, before revising any topic lesson.}
\end{frame}

\begin{frame}{The Four Depth Tags}
  \begin{itemize}
    \item \key{Core-General}: stable fundamentals every post tests (bit is smallest unit; F5 starts slide show; RAM is volatile).
    \item \key{Current-General}: stable concept whose current form must be re-checked (latest anti-virus; role of IT in governance; portal facts).
    \item \key{Technical-Post}: deeper material asked only of Website Operator / Computer Instructor (cache levels; UTF-8 variable length; binary conversion drills).
    \item \key{Version-Sensitive}: any claim that depends on a software or protocol version (Office shortcut; POP3 behaviour; HTTP version transport).
  \end{itemize}
  \begin{block}{Rule}
    Every item gets exactly one depth tag plus a version flag where needed.
  \end{block}
\end{frame}

\begin{frame}{How to Tag Any Item}
  \begin{itemize}
    \item Ask 1: does the \key{Junior Assistant scope} (Unit IV) name this concept? If yes, start at \key{Core-General}.
    \item Ask 2: does the answer change with time, portal, or version? If yes, add \key{Current-General} or \key{Version-Sensitive}.
    \item Ask 3: does it need engineering depth (multi-step conversion, protocol internals)? If yes, mark \key{Technical-Post}.
    \item Worked tagging: ``F5 starts slide show'' $\rightarrow$ \key{Core-General}; ``which menu holds this Word option'' $\rightarrow$ \key{Version-Sensitive}.
  \end{itemize}
  \begin{exampleblock}{Worked tagging}
    ``IPv6 is 128-bit'' is Core-General; ``this protocol version rides on QUIC'' is Version-Sensitive; ``convert 11010101.01'' is Technical-Post.
  \end{exampleblock}
\end{frame}

\begin{frame}{The Official Source Ledger (Tier 1)}
  \begin{itemize}
    \item \key{SRC-01:} Notification 08 of 2025 with the official JA scope (Unit IV) and the 10-mark annexure.
    \item \key{SRC-02:} JKSSB question-paper index: which papers exist per year and post.
    \item \key{SRC-03:} Junior Assistant 2026 paper, computer block Q61--80.
    \item \key{SRC-04:} Website Operator 2026 paper, computer block circa Q61--80.
  \end{itemize}
  \muted{The ledger continues on the next frame; Tier 1 always outranks banks and textbooks.}
\end{frame}

\begin{frame}{The Official Source Ledger (continued)}
  \begin{itemize}
    \item \key{SRC-05:} Computer Instructor-Operator 2025 paper, 120 items.
    \item \key{SRC-06:} JA 2026 final answer key: settles disputed JA items.
    \item \key{SRC-07:} WO 2026 provisional key: provisional only, mark its status.
    \item \key{SRC-08:} CI 2025 final key: settles disputed technical-post items.
  \end{itemize}
  \begin{alertblock}{Authority rule}
    A Tier-1 final key \key{supersedes} every bank, textbook, and tutorial on a disputed item.
  \end{alertblock}
\end{frame}

\begin{frame}{Lower Tiers: What They Are For}
  \begin{itemize}
    \item \key{Tier 2 (banks):} Sanfoundry and standard MCQ compilations supply practice items and distractors only.
    \item Tier-2 answers are \key{never authoritative}; re-verify each one and label it adapted, never PYQ.
    \item \key{Tier 3 (reference):} vendor docs, portal pages, textbooks explain concepts; they settle no exam dispute.
    \item Vendor docs must be tied to the \key{prescribed Office version} before a shortcut claim is trusted.
  \end{itemize}
  \begin{block}{Keep the tiers straight}
    Practice from Tier 2; verify against Tier 1; explain with Tier 3.
  \end{block}
\end{frame}

\begin{frame}{How to Check Answer-Key Changes}
  \begin{itemize}
    \item Place the \key{provisional key} and the \key{final key} side by side for your post.
    \item Mark every item where the answer \key{changed}, was dropped, or gained an alternate answer.
    \item Changed items are \key{defective-option signals}: quote the option verbatim, flag the defect, teach the stable concept.
    \item Never memorise a provisional answer as final; re-check before the next mock.
  \end{itemize}
  \begin{exampleblock}{Worked check}
    If a provisional key later changes on a disputed wording item, the final key wins and the stem wording becomes a flagged example, not a fact.
  \end{exampleblock}
\end{frame}

\begin{frame}{Why Technical Depth Must NOT Be Generalised}
  \begin{itemize}
    \item The Junior Assistant paper tests \key{recall and classification}; the Computer Instructor paper tests \key{depth and computation}.
    \item Example: binary-to-decimal conversion drills and encoding details belong to the \key{Technical-Post} tag.
    \item A general-post candidate who masters the whole technical block \key{wastes time} and risks L4/L5 distractors.
    \item A technical-post candidate who studies only the general block \key{under-prepares} for the deeper items.
  \end{itemize}
  \begin{alertblock}{Anti-pattern}
    Generalising Website Operator / Computer Instructor depth to every post over-scopes the general candidate.
  \end{alertblock}
\end{frame}

\begin{frame}{Two Concrete Non-Generalisations}
  \begin{itemize}
    \item \key{Foundational JA item:} ``smallest unit of data is the bit'' --- Core-General, every post.
    \item \key{Technical CI item:} ``binary 11010101.01 equals 213.25'' --- Technical-Post, computation drill.
    \item Studying the second item's method deeply is \key{required} for CI, \key{optional} for JA.
    \item The routing rule: JA masters the first; CI masters both; neither confuses one's depth for the other's.
  \end{itemize}
  \muted{Tag first, then allocate study time by tag, not by enthusiasm.}
\end{frame}

\begin{frame}{Version Sensitivity: Office}
  \begin{itemize}
    \item Office claims are \key{version-sensitive}: menus, shortcuts, and SmartArt behaviour move between versions.
    \item Always tie an Office fact to the \key{prescribed version} via vendor documentation.
    \item ``Which tab holds this command'' is answerable only for a \key{stated version}; without one, flag the item.
    \item Stable Office facts (Mail Merge merges letters with a database; formulas start with \key{=}) stay Core-General.
  \end{itemize}
  \begin{block}{Rule}
    Version-stated Office facts can be trusted; version-free menu claims cannot.
  \end{block}
\end{frame}

\begin{frame}{Version Sensitivity: Protocols and Keys}
  \begin{itemize}
    \item Protocol claims are \key{version-sensitive}: POP3 downloads locally, IMAP syncs; HTTP/3 uses QUIC over UDP, not TCP.
    \item A protocol item is settled by the \key{final key plus the standard}, never by a bank alone.
    \item IPv6 is \key{128-bit} with a 40-byte header; a 256-bit option is a planted distractor.
    \item Digital signature means \key{authenticity and integrity}; encryption means confidentiality --- versions do not change this split.
  \end{itemize}
  \begin{alertblock}{Trap alert}
    ``POP3 keeps mail synced on the server'' and ``IPv6 is 256-bit'' are version-confusion distractors.
  \end{alertblock}
\end{frame}

\begin{frame}{Trap Alert: Audit-Stage Errors}
  \begin{itemize}
    \item Treating a \key{provisional key} as final and memorising its answers.
    \item Presenting a generated item as a \key{real PYQ} instead of labelling it Original.
    \item Importing \key{engineering depth} into a general-post revision set.
    \item Quoting a \key{defective option} as fact instead of flagging it and teaching the stable concept.
    \item Citing a Tier-2 bank to \key{overrule} a Tier-1 final key.
  \end{itemize}
\end{frame}

\begin{frame}{A One-Page Audit Checklist}
  \begin{itemize}
    \item New notification read; scope confirmations noted (SRC-01).
    \item Newest paper collected (SRC-03/04/05); provisional and final keys collected (SRC-06/07/08).
    \item Every computer item tagged: \key{Core-General / Current-General / Technical-Post / Version-Sensitive}.
    \item Key changes diffed; defective options flagged; version-sensitive claims tied to a version.
    \item Study time routed by post; banks re-labelled; no generated item called a PYQ.
  \end{itemize}
  \muted{File this checklist with your notes each cycle; it is the audit trail.}
\end{frame}

\begin{frame}{Lesson 66: Recall}
  \begin{itemize}
    \item The five-step workflow: collect Tier 1, tag each item, diff provisional vs final key, flag versions, route by post.
    \item Four tags: \key{Core-General}, Current-General, Technical-Post, Version-Sensitive.
    \item Ledger: notification and scope (SRC-01), paper index (SRC-02), papers (SRC-03/04/05), keys (SRC-06/07/08).
    \item Final keys settle disputes; provisional keys are marked provisional; banks never overrule Tier 1.
    \item Technical-post depth is \key{not generalised}; Office and protocol claims are version-checked.
  \end{itemize}
\end{frame}
\end{document}
